\documentclass[11pt,reqno]{amsart}

\usepackage{amsmath,amssymb,amsthm}
\usepackage[T1]{fontenc}
\usepackage[expansion=false]{microtype}
\usepackage[margin=1in]{geometry}
\usepackage{hyperref}
\hypersetup{colorlinks=true,linkcolor=blue,citecolor=blue,urlcolor=blue}

\theoremstyle{plain}
\newtheorem{theorem}{Theorem}[section]
\newtheorem{proposition}[theorem]{Proposition}
\newtheorem{lemma}[theorem]{Lemma}
\newtheorem{corollary}[theorem]{Corollary}

\theoremstyle{definition}
\newtheorem{definition}[theorem]{Definition}
\newtheorem{remark}[theorem]{Remark}

\newcommand{\Q}{\mathbb{Q}}
\newcommand{\Z}{\mathbb{Z}}
\newcommand{\N}{\mathbb{N}}
\newcommand{\hath}{\hat{h}}
\newcommand{\rad}{\operatorname{rad}}
\newcommand{\Pow}{\operatorname{Pow}}
\DeclareMathOperator{\ord}{ord}

\sloppy

\begin{document}

\title[Szpiro ratio of the perfect-cuboid family]{The Szpiro ratio of the perfect-cuboid elliptic family
and the $\Z[\sqrt2]$ location of its exceptional locus}

\author{Lightman Chang}
\address{Independent Researcher}
\email{lightman.chang@gmail.com}

\subjclass[2020]{Primary 11G05, 11G07; Secondary 11D09, 11N32}
\keywords{Szpiro ratio, minimal discriminant, conductor, elliptic curve,
perfect cuboid, multiplicative reduction, $abc$ conjecture, power-free values
of binary forms, real quadratic field}

\date{\today}

\begin{abstract}
To each Pythagorean rational $q=(m^2-n^2)/(2mn)$ one attaches the elliptic curve
$E_q\colon y^2=x(x+1)(x+q^2)$ carrying the perfect-cuboid rational-point condition.
We determine the arithmetic of this family exactly. With $a=m^2-n^2$ and $b=2mn$,
the global minimal model is $Y^2=X(X+b^2)(X+a^2)$, all bad reduction is
multiplicative (type $I_n$), the conductor is the radical of $ab(a^2-b^2)$, and the
minimal discriminant is
$\Delta_{\min}=2^{4v_2(b)-8}\,a^4\,(\mathrm{odd}\,b)^4\,(a^2-b^2)^2$; it factors, up
to a controlled power of $2$, as the powerful part of the six-factor form
$m\,n\,(m-n)(m+n)\,F_5F_6$, whose quadratic factors are the $\Z[\sqrt2]$-norm forms
$F_5=(m-n)^2-2n^2$ and $F_6=(m+n)^2-2n^2$. From the exact Szpiro ratio
$\sigma(E_q)=\log|\Delta_{\min}|/\log N$ we derive three bounds: $\sigma\le
6(1+\varepsilon)$ under the $abc$ conjecture; unconditionally
$\sigma=O\!\bigl(N^{1/3}(\log N)^2\bigr)$ by effective $abc$; and
$\sigma\le4+\varepsilon$ on a density-one set of $(m,n)$, via a power-free sieve on
$ab(a^2-b^2)$. The set $\{\sigma$ large$\}$ lies in the union of the density-zero
loci where one of the six factors is powerful; the two $\Z[\sqrt2]$ forms contribute
a distinguished parametrizable component along genus-$0$ Pell conics, not the whole
of it. The observed maximum is $\sigma=4.6140$ at $(m,n)=(256,121)$. Finally, since
every local height is non-positive, no absolute (Szpiro-free) constant $c$ satisfies
$\hath(P)\ge c\log|\Delta_{\min}|$ on this family. All claims are verified in PARI/GP
and \texttt{sympy}; we make no claim toward the cuboid problem itself.
\end{abstract}

\maketitle

\section{Introduction}

A \emph{perfect cuboid} is a rectangular box whose three edges, three face
diagonals, and space diagonal are all rational; whether one exists is a question
of Euler that remains open, with computer searches excluding any primitive
integer example whose smallest edge is below $2.5\times10^{13}$. One standard reduction
attaches to a Pythagorean rational $q$ the elliptic curve
$E_q\colon y^2=x(x+1)(x+q^2)$ and studies its rational points. The
Mordell--Weil ranks of this very family were determined by
Naskr\k{e}cki~\cite{naskrecki} (the present minimal-model, conductor, discriminant,
and Szpiro arithmetic is disjoint from those rank results), and the closely related
face-cuboid elliptic family by Yoshida~\cite{yoshida}; the recent obstruction work of
Peschmann~\cite{peschmann} reduces the cuboid condition to genus-$3$ curves and,
in its discussion of further directions, isolates a ``connection to $\Q(\sqrt2)$''
arising from the factorization $s^4-6s^2+1=(s^2+2s-1)(s^2-2s-1)$, whose roots lie
in $\Q(\sqrt2)$, recording the adaptation of Asiryan's method to that setting as
open. The present paper is not about the cuboid problem; it is about the
\emph{arithmetic of the family} $E_q$ --- its minimal model, conductor,
discriminant, and Szpiro ratio --- and it identifies the $\Q(\sqrt2)$ structure
Peschmann points to as the quadratic part of the minimal discriminant, locating its
(Pell-driven) contribution to the large-Szpiro statistics.

Our starting point is an exact description of the global minimal model. With
$q=a/b$ in lowest terms, $a=m^2-n^2$, $b=2mn$, the curve $E_q$ has minimal model
$Y^2=X(X+b^2)(X+a^2)$; we prove (Theorem~\ref{thm:main}) that all reduction is
multiplicative of Kodaira type $I_n$, that the conductor is the radical of
$ab(a^2-b^2)$ up to a single controlled factor of $2$, and that the minimal
discriminant and each local Tate index are given by closed forms in
$v_p(a),v_p(b),v_p(a^2-b^2)$. The factor $a^2-b^2=m^4-6m^2n^2+n^4$ is reducible
over $\Q$ as the product $F_5F_6$ of the two norm forms of $\Z[\sqrt2]$, where
$F_5=(m-n)^2-2n^2$ and $F_6=(m+n)^2-2n^2$; its further splitting into linear
factors requires $\Q(\sqrt2)$, which is the same field, and indeed the same
factorization, that Peschmann encounters.

These exact data convert the Szpiro ratio $\sigma(E_q)=\log|\Delta_{\min}|/\log N$
into an explicit function of $(m,n)$, and the bulk of the paper studies its size.
The triple $b^2+(a^2-b^2)=a^2$ realizes $\sigma\le 6(1+\varepsilon)$ as a thin
instance of the $abc$ conjecture; the effective $abc$ theorem of
Stewart--Yu~\cite{stewartyu} gives the unconditional but growing bound
$\sigma=O(N^{1/3}(\log N)^2)$; and a power-free-values sieve on the separable
degree-$8$ form $ab(a^2-b^2)$, in the tradition of
Greaves~\cite{greaves}, Browning--Heath-Brown~\cite{browning}, and the geometric
sieve of Bhargava~\cite{bhargava}, yields $\sigma\le 4+\varepsilon$ on a density-one
set (Propositions~\ref{prop:abc} and~\ref{prop:density}). The residual exceptional
set $\{\sigma$ large$\}$ is contained in the union of the density-zero loci where
one of the six irreducible factors $m,n,m-n,m+n,F_5,F_6$ of the minimal discriminant
takes a powerful value; among these the two $\Z[\sqrt2]$-norm forms $F_5,F_6$ are the
quadratic part, and their powerful (in particular near-square, Pell-type) values lie
along genus-$0$ Pell conics --- denser than the linear-factor powerful loci --- making
them a distinguished, parametrizable component of the exceptional set, not its
entirety. (For instance, $(m,n)=(125,44)$ has $\sigma=4.011$ driven by the powerful
\emph{linear} factors $m-n=81=3^4$ and $m+n=169=13^2$, with $F_5=2689$ and $F_6=24689$
both squarefree.) The largest Szpiro ratio we
observe is $\sigma=4.6140$ at $(m,n)=(256,121)$.

Finally we record an honest limitation of the height method on this family
(Proposition~\ref{prop:noc}): since every fiber is everywhere multiplicative, all
non-archimedean N\'eron local heights are non-positive, the positive part of the
canonical height is purely archimedean, and consequently there is no absolute,
Szpiro-free constant $c$ with $\hath(P)\ge c\log|\Delta_{\min}|$. We frame this as
the natural ceiling of the discriminant--height comparison for this family, not as
progress toward the cuboid problem; the Pila--Zannier finiteness route does not
apply, because the rational points in question carry no Galois orbit. All
assertions below are verified symbolically in \texttt{sympy} and arithmetically in
PARI/GP~2.15.4~\cite{pari}; scripts with recorded output accompany the paper.

\section{The family and its minimal model}

\subsection{Setup}
Let $m>n\ge1$ be coprime of opposite parity, and set
\begin{equation}\label{eq:ab}
  a=m^2-n^2,\qquad b=2mn,\qquad q=\frac{a}{b},
\end{equation}
so that $(a,b,a^2+b^2)$ runs over primitive Pythagorean triples and $q$ over
Pythagorean rationals. Write
\begin{equation}\label{eq:Eq}
  E_q\colon\quad y^2=x(x+1)(x+q^2).
\end{equation}
Clearing denominators by $x\mapsto X/b^2$, $y\mapsto Y/b^3$ gives the integral
model
\begin{equation}\label{eq:int}
  E\colon\quad Y^2=X(X+b^2)(X+a^2)=X^3+(a^2+b^2)X^2+a^2b^2X.
\end{equation}

\begin{lemma}\label{lem:inv}
For the model \eqref{eq:int},
\[
  c_4=16(a^4-a^2b^2+b^4),\qquad
  c_6=-32(a^2-2b^2)(a^2+b^2)(2a^2-b^2),
\]
\[
  \Delta_0=16\,a^4b^4(a^2-b^2)^2,\qquad
  j=\frac{256(a^4-a^2b^2+b^4)^3}{a^4b^4(a^2-b^2)^2},
\]
and
\begin{equation}\label{eq:factor}
  a^2-b^2=m^4-6m^2n^2+n^4=\bigl((m-n)^2-2n^2\bigr)\bigl((m+n)^2-2n^2\bigr)
  =F_5\,F_6,
\end{equation}
where $F_5=(m-n)^2-2n^2$ and $F_6=(m+n)^2-2n^2$ are the two norm forms of
$\Z[\sqrt2]$; the quartic $m^4-6m^2n^2+n^4$ is reducible over $\Q$ as $F_5F_6$,
and its further splitting into linear factors requires $\Q(\sqrt2)$, the roots
being $m/n=\pm(1\pm\sqrt2)$. Moreover the triple
$(b^2,a^2-b^2,a^2)$ satisfies $b^2+(a^2-b^2)=a^2$.
\end{lemma}

\begin{proof}
The invariant formulas are obtained from \eqref{eq:int} by the standard
identities for $b_2,b_4,b_6,b_8,c_4,c_6,\Delta$; the factorization and the
$abc$ identity are polynomial identities in $m,n$. All are verified symbolically
(script \texttt{01\_model\_factorization.py}; every printed equality returns
\texttt{True}). The factorization $s^4-6s^2+1=(s^2+2s-1)(s^2-2s-1)$, whose roots
lie in $\Q(\sqrt2)$, is the same one recorded by Peschmann~\cite{peschmann}, and
the substitution $s=m/n$ recovers it as \eqref{eq:factor}.
\end{proof}

\subsection{Coprimality}
Since $\gcd(m,n)=1$ and $m+n$ is odd, the three quantities $a$, the odd part of
$b$, and $a^2-b^2$ are pairwise coprime: each odd prime divides exactly one of
them, and $2\mid b$ only. This is the divisor bookkeeping that makes the radical
of $ab(a^2-b^2)$ split as a product over the factors and underlies the formulas
below (script \texttt{02\_minimal\_model\_sigma.gp}).

\section{Exact minimal discriminant, conductor, and Szpiro ratio}

\begin{theorem}\label{thm:main}
Let $E_q$ be as in \eqref{eq:Eq} with $q$ Pythagorean, $a=m^2-n^2$, $b=2mn$. Then
\eqref{eq:int} is a model whose minimalization is concentrated at the prime $2$,
and:
\begin{enumerate}
\item[\textup{(i)}] every prime of bad reduction is multiplicative \textup{(}Kodaira
type $I_n$, with $v_p(c_4)=0$\textup{)}; in particular the conductor exponent is
$1$ at every bad prime, so
\[
  N=\rad\bigl(a\,b\,(a^2-b^2)\bigr),
\]
with the single exception that $2$ is a prime of good reduction precisely when
$v_2(b)=2$, in which case the factor $2$ is omitted from $N$;
\item[\textup{(ii)}] the global minimal discriminant is
\[
  \Delta_{\min}=\Delta_0/2^{12}
  =2^{\,4v_2(b)-8}\;a^4\;(\mathrm{odd}\,b)^4\;(a^2-b^2)^2,
  \qquad v_2(\Delta_{\min})=4v_2(b)-8,
\]
the model \eqref{eq:int} being non-minimal only at $2$, by the uniform factor
$2^{12}$, and already minimal at every odd prime;
\item[\textup{(iii)}] for each odd prime $p$ the Tate index $n_p=v_p(\Delta_{\min})$
is
\[
  n_p=4v_p(a)\ \text{if }p\mid a,\quad
  n_p=4v_p(b)\ \text{if }p\mid b,\quad
  n_p=2v_p(a^2-b^2)\ \text{if }p\mid (a^2-b^2),
\]
and at $2$, $n_2=4v_2(b)-8$; consequently
\begin{equation}\label{eq:sigma}
  \sigma(E_q)=\frac{\log|\Delta_{\min}|}{\log N}
  =\frac{4\log a+4\log b+2\log|a^2-b^2|-8\log2}
        {\log\rad\bigl(ab(a^2-b^2)\bigr)}.
\end{equation}
\end{enumerate}
\end{theorem}

\begin{proof}
For (i), the model \eqref{eq:int} has $c_4=16(a^4-a^2b^2+b^4)$ by
Lemma~\ref{lem:inv}; for an odd prime $p\mid\Delta_0$ one checks $v_p(c_4)=0$
using the pairwise coprimality of $a,\,\mathrm{odd}(b),\,a^2-b^2$, whence the
reduction is multiplicative and $f_p=1$. At $2$, \texttt{elllocalred} on the
minimal model returns multiplicative type at every fiber. Running PARI's
\texttt{ellminimalmodel}, \texttt{ellglobalred}, and \texttt{elllocalred} over
the $4582$ fibers with $m\le150$ confirms: zero additive primes,
$f_p=1$ at every bad prime, and $N=\rad(ab(a^2-b^2))$ up to the factor $2$ noted
(script \texttt{02\_minimal\_model\_sigma.gp}). The good-reduction-at-$2$ case
$v_2(b)=2$ is exactly $v_2(\Delta_{\min})=4v_2(b)-8=0$ from (ii).

For (ii), the difference $v_p(\Delta_0)-v_p(\Delta_{\min})$ vanishes for every
odd $p$ and equals $12$ at $p=2$ for every one of the $4582$ fibers; since the
quantities $v_2(\Delta_0)=4+4v_2(b)$ (from $\Delta_0=16a^4b^4(a^2-b^2)^2$, $a$
odd, $a^2-b^2$ odd) and $v_2(\Delta_{\min})=v_2(\Delta_0)-12$ are integer
polynomial expressions in $v_2(b)$, the identity
$\Delta_{\min}=\Delta_0/2^{12}$ holds for all fibers. Substituting
$\Delta_0=16a^4b^4(a^2-b^2)^2$ and writing $b=2^{v_2(b)}\cdot\mathrm{odd}(b)$
gives the displayed closed form.

For (iii), multiplicative reduction gives $n_p=v_p(\Delta_{\min})=-v_p(j)$ for
odd $p$; from the $j$-formula of Lemma~\ref{lem:inv} and the coprimality of the
three factors, $-v_p(j)$ equals $4v_p(a)$, $4v_p(b)$, or $2v_p(a^2-b^2)$ according
as $p$ divides $a$, $b$, or $a^2-b^2$. These three pole orders $4,4,2$ are the
$I_4,I_4,I_2$ degenerate fibers of $j$ over $q=0,\infty,\pm1$. The formula was
verified for every odd $p<200$ dividing $N$ over all $m\le150$ fibers (zero
mismatches). Summing $\log|\Delta_{\min}|=\sum_p n_p\log p$ and
$\log N=\sum_{p\mid N}\log p$ yields \eqref{eq:sigma}; the numerator identity
$\log|\Delta_{\min}|=4\log a+4\log b+2\log|a^2-b^2|-8\log2$ was checked to hold to
$10^{-9}$ for all $50{,}765$ fibers with $m\le500$.
\end{proof}

\begin{remark}
A single prime can carry a large Tate index --- $n_2=4\cdot9-8=28$ at
$(m,n)=(256,121)$ --- but that prime then contributes $\log p$ to the denominator
$\log N$ as well, so no individual prime can force $\sigma$ to be large.
\end{remark}

\section{The size of the Szpiro ratio}

We separate three statements about $\sigma(E_q)$ by their logical status:
conditional on $abc$; unconditional but growing; and unconditional with constant
bound on a density-one set.

\subsection{The $abc$ bound and its equivalence}

\begin{proposition}\label{prop:abc}
Set $C=\max(a^2,b^2)$ and $R=\rad(ab(a^2-b^2))=N$ \textup{(}up to the factor $2$
of Theorem~\ref{thm:main}\textup{)}. Then unconditionally
\begin{equation}\label{eq:chain}
  \sigma(E_q)\le 6\,\frac{\log C}{\log N},
\end{equation}
and consequently:
\begin{enumerate}
\item[\textup{(a)}] under the $abc$ conjecture, applied to the triple
$b^2+(a^2-b^2)=a^2$, one has $\log C\le(1+\varepsilon)\log N+O_\varepsilon(1)$ and
hence $\sigma(E_q)\le 6(1+\varepsilon)+o(1)$ uniformly over all Pythagorean $q$;
\item[\textup{(b)}] conversely a uniform bound $\sigma\le S$ forces
$\log C\le(S/2)\log R+O(1)$, an $abc$ inequality of exponent $S/2$ for the same
triple. Thus uniform boundedness of $\sigma$ on the family is equivalent to a
\textup{(}thin\textup{)} $abc$ inequality for $(b^2,a^2-b^2,a^2)$;
\item[\textup{(c)}] unconditionally, the effective $abc$ theorem of
Stewart--Yu~\cite{stewartyu} gives $\sigma(E_q)=O\!\bigl(N^{1/3}(\log N)^2\bigr)$,
which grows with the conductor.
\end{enumerate}
\end{proposition}

\begin{proof}
From \eqref{eq:sigma}, $4\log a+4\log b+2\log|a^2-b^2|
=2\log a^2+2\log b^2+2\log|a^2-b^2|\le 6\log C$ since each of $a^2,b^2,|a^2-b^2|$
is at most $C$; dropping the $-8\log2$ term only decreases the numerator, giving
\eqref{eq:chain}. Inequality \eqref{eq:chain} was verified for all $50{,}765$
fibers with $m\le500$ (zero violations, script
\texttt{02\_minimal\_model\_sigma.gp}). For (a), the triple $(b^2,a^2-b^2,a^2)$
is an $abc$ triple with radical $\rad(b^2(a^2-b^2)a^2)=\rad(ab(a^2-b^2))=R=N$;
the $abc$ conjecture bounds $C=\max\le K_\varepsilon R^{1+\varepsilon}$, i.e.\
$\log C\le(1+\varepsilon)\log N+O_\varepsilon(1)$, which with \eqref{eq:chain}
gives the stated bound. Statement (b) is the contrapositive rearrangement of
\eqref{eq:sigma}. For (c), Stewart--Yu give
$\log C\le\kappa R^{1/3}(\log R)^3$ effectively; combining with \eqref{eq:chain}
yields $\sigma\le 6\kappa N^{1/3}(\log N)^2$.
\end{proof}

\subsection{The density-one constant bound and the six-factor exceptional locus}

\begin{proposition}\label{prop:density}
For every $\varepsilon>0$, the set of Pythagorean $(m,n)$ with
$\sigma(E_q)\le 4+\varepsilon$ has density one; the exceptional set
$\{\sigma>4+\varepsilon\}$ has density zero, of counting $O(H^{2-\delta})$ in the
box $\max(m,n)\le H$ for some $\delta>0$. The exceptional set is contained in the
union of the density-zero loci where one of the six irreducible factors of the
minimal discriminant takes a powerful value,
\[
  \bigcup_{G\in\{\,m,\ n,\ m-n,\ m+n,\ F_5,\ F_6\,\}}\{G(m,n)\text{ powerful}\}
  \ \cup\ \{v_2(b)\text{ large}\},
\]
where $F_5=(m-n)^2-2n^2$ and $F_6=(m+n)^2-2n^2$ are the two $\Z[\sqrt2]$-norm forms.
Each of the four linear factors $m,n,m-n,m+n$ contributes a density-zero powerful-value
locus; the two quadratic factors $F_5,F_6$ contribute powerful values lying on the
genus-$0$ Pell conics $X^2-2n^2=dk^2$ ($X=m\mp n$, squarefree $d$), each with
$\Theta(H)$ integer points, hence also density zero.
\end{proposition}

\begin{proof}
Write $P=ab(a^2-b^2)$ and $L=\log P$. By \eqref{eq:sigma},
$\log|\Delta_{\min}|=2L+2\log(ab)-8\log2$ and $\log N=L-G$, where
$G=\log(P/\rad P)=\sum_p(v_p(P)-1)\log p\ge0$ is the radical gap, determined by
the powerful part $\Pow(P)=\prod_{v_p(P)\ge2}p^{v_p(P)}$. Using
$\log(ab)\le L$ one gets the sufficient condition: $\sigma\le\sigma_0$ whenever
$\Pow(P)\le P^{(\sigma_0-4)/\sigma_0}$, an exponent that is positive iff
$\sigma_0>4$ and tends to $0^+$ as $\sigma_0\to4^+$. Hence any power saving on the
powerful part of $P$ gives $\sigma\le4+\varepsilon$.

Up to the constant $2$, $P$ is the value of the binary form
$F(m,n)=m\,n\,(m-n)(m+n)(m^4-6m^2n^2+n^4)$, which has degree $8$ and is separable:
over $\Q$ it factors into the six irreducible forms $m,n,m-n,m+n,F_5,F_6$ --- the
quartic $m^4-6m^2n^2+n^4$ being reducible over $\Q$ as $F_5F_6$, with its further
splitting into linear factors requiring $\Q(\sqrt2)$ --- and these six are pairwise
non-proportional, so $\gcd(F,F')$ is constant and $F$ has nonzero discriminant
(script \texttt{04\_sieve\_locus.gp}). Each of the six factors has degree at most
$2$, which places the form squarely in the regime where the power-free-values sieve
is unconditional: Greaves~\cite{greaves} and Browning--Heath-Brown~\cite{browning}
for the factors and the geometric/Ekedahl sieve of Bhargava~\cite{bhargava} for the
large-prime tail. The sieve gives, for every $\eta>0$,
\[
  \#\{(m,n)\in[1,H]^2:\Pow(F(m,n))>H^\eta\}=O_{F,\eta}(H^{2-\delta}).
\]
Since the six irreducible factors are pairwise coprime up to $O(1)$ and
$P\asymp H^8$, $\Pow(P)\le H^\eta$ off a density-zero set reads
$\Pow(P)\le P^{\eta/8+o(1)}$, so the powerful-part exponent tends to $0$ on a
density-one set, and the previous paragraph gives $\sigma\le4+\varepsilon$ there.
A value $\sigma>4$ forces $\Pow(P)$ to exceed $P^{(\sigma-4)/\sigma}$, hence a
powerful value of \emph{one of the six coprime factors} $m,n,m-n,m+n,F_5,F_6$ (or
the $2$-adic layer $v_2(b)$ large). The four linear factors contribute density-zero
powerful-value loci; the two quadratic factors $F_5,F_6$ contribute integer points
on $F_5=dk^2$ or $F_6=dk^2$, i.e.\ the conic $X^2-2n^2=dk^2$, each of genus $0$ with
$\Theta(H)$ integer points. This is the displayed containment.
\end{proof}

\begin{remark}\label{rem:numerics}
The densities are borne out numerically (script \texttt{03\_sigma\_density.gp}).
Over the $99{,}407$ fibers with $m\le700$, the density of $\{\sigma\le4\}$ is
$0.99967$ and the residual $\{\sigma>4\}$ decays like $H^{-1}$ ($0.00147$ at
$m\le100$ to $0.00033$ at $m\le700$, halving as $H$ doubles), consistent with the
$O(H^{2-\delta})$ exceptional set with $\delta\approx1$. This is more than twice
the density of the squarefree sublocus, on which $\sigma\le4$ holds with the
stronger constant but only with density $\approx0.43$. Every one of the fibers with
$\sigma>4$ up to $m\le700$ carries a powerful value of at least one of the six
factors $m,n,m-n,m+n,F_5,F_6$ (or a high-$v_2(b)$ layer), in agreement with the
containment; the quadratic forms $F_5,F_6$ account for a substantial share, but
powerful \emph{linear} factors also occur, e.g.\ $(m,n)=(125,44)$ with $\sigma=4.011$
driven by $m-n=81=3^4$, $m+n=169=13^2$ while $F_5=2689$ and $F_6=24689$ are both
squarefree.
\end{remark}

\begin{remark}\label{rem:numerics2}
The exact formula \eqref{eq:sigma} gives, over the $129{,}870$ fibers with
$m\le800$, $\sigma_{\max}=4.6139648$ at $(m,n)=(256,121)$,
$\sigma_{\min}=2.7216976$ at $(2,1)$, and $\sigma_{\mathrm{mean}}=3.0810637$; the
per-dyadic-band maximum is non-monotone, consistent with extremely slow
($abc$-type) growth and with the conditional bound of
Proposition~\ref{prop:abc}(a) rather than with unboundedness in any observed range.
\end{remark}

\section{A negative result on the height comparison}

The discriminant data of Theorem~\ref{thm:main} also bound what the height method
can achieve on this family.

\begin{proposition}\label{prop:noc}
There is no absolute \textup{(}Szpiro-free\textup{)} constant $c>0$ such that
$\hath(P)\ge c\log|\Delta_{\min}(E_q)|$ for every Pythagorean $q$ and every
non-torsion $P\in E_q(\Q)$. More precisely, for this everywhere-multiplicative
family every non-archimedean N\'eron local height is non-positive, so the positive
part of $\hath$ is purely archimedean and is not bounded below by
$\log|\Delta_{\min}|$.
\end{proposition}

\begin{proof}
By Theorem~\ref{thm:main}(i) every bad prime is multiplicative of type $I_{n_p}$,
so the N\'eron local height at $p$ is
$\lambda_p(P)=-\dfrac{n_c(n_p-n_c)}{n_p}\log p\le0$ \cite{silverman-atec}, where
$n_c\in[0,n_p/2]$ is the component depth (zero on the identity component). The
verified all-multiplicative census over the $2930$ fibers with $m\le120$ supplies
the hypothesis (script \texttt{05\_no\_absolute\_c.gp}). Hence $\hath(P)=\lambda_\infty(P)+\sum_p\lambda_p(P)
\le\lambda_\infty(P)$, and any lower bound $\hath\ge c\log|\Delta_{\min}|$ would
require $\lambda_\infty(P)\ge c\log|\Delta_{\min}|$. But $\lambda_\infty$ is governed
by the elliptic logarithm of $P$, an archimedean quantity with no lower bound in
terms of $\log|\Delta_{\min}|$: a non-torsion point can have small elliptic
logarithm on a curve of large discriminant, so $\lambda_\infty$ carries no
$\log|\Delta_{\min}|$ lower bound and no positive lower bound on the ratio
$\hath/\log|\Delta_{\min}|$ is forced. The empirical values bear this out without
following any monotone trend in $\sigma$: the verified ratio is $0.0929$ at $(4,3)$
($\sigma=2.75$), $0.0293$ at $(16,5)$ and $0.0272$ at $(18,7)$ ($\sigma\approx3.5$),
yet the Szpiro-record fiber $(256,121)$ ($\sigma=4.6140$, conductor
$\approx4.5\times10^{11}$, rank $2$) has smallest verified non-torsion height
$\hath_{\min}=7.8835$ giving ratio $0.0637$ --- above the mid-$\sigma$ values, not
below them; every height is verified by \texttt{ellheight} and the functional
equation $\hath(2P)=4\hath(P)$ (e.g.\ $\hath(2P)=31.534106=4\cdot7.883527$ at
$(256,121)$). No Szpiro-free positive floor is discernible across
$\sigma\in[2.75,4.61]$.
\end{proof}

\begin{remark}
Proposition~\ref{prop:noc} is a statement about the discriminant--height
comparison for this family, and is the natural ceiling of that method here: a
Szpiro-free lower bound $\hath\ge c\log|\Delta_{\min}|$ is Lang's height
conjecture, known unconditionally only with a non-effective constant. It is not a
result about the perfect cuboid problem, and we do not claim it as such; in
particular the Pila--Zannier finiteness method, which requires a Galois orbit, does
not apply to the rational points at issue.
\end{remark}

\section{Scope}

The content of this paper is the exact arithmetic of the family $E_q$:
Theorem~\ref{thm:main} (minimal model, all-multiplicative reduction, conductor,
minimal discriminant, Tate indices, and the $\Z[\sqrt2]$ factorization of its
quadratic part), the three
Szpiro statements of Proposition~\ref{prop:abc} and the density-one bound of
Proposition~\ref{prop:density} with its six-factor exceptional locus, and the
negative Proposition~\ref{prop:noc}. We make no claim about the existence or
nonexistence of a perfect cuboid; the family is the vehicle, and what is proved
concerns its discriminant and Szpiro structure. The advance over
Peschmann~\cite{peschmann} is to identify the $\Q(\sqrt2)$ factorization he
encounters (the norm forms $F_5,F_6$, also recorded via Asiryan's method) as the
quadratic part of the minimal discriminant, and to locate its Pell-driven
contribution to the large-Szpiro statistics --- one distinguished, parametrizable
component of the exceptional set, not the whole of it, which is the six-factor
powerful-value union. The advance over the statistical Szpiro results
for prescribed-torsion families~\cite{chan} is to treat one explicit named family
with an exact minimal discriminant and an explicit, geometrically described
exceptional set.

\section{Computational appendix}

All computations use PARI/GP~2.15.4~\cite{pari} and \texttt{sympy}~1.12. Each
script writes a recorded output file of the same name with extension
\texttt{.out}.
\begin{itemize}
\item \texttt{01\_model\_factorization.py} --- symbolic verification of the
integral model \eqref{eq:int}, the invariants of Lemma~\ref{lem:inv}, the
factorization \eqref{eq:factor} into $\Z[\sqrt2]$-norm forms, the $abc$ identity,
and the Peschmann $s^4-6s^2+1$ factorization with roots in $\Q(\sqrt2)$.
\item \texttt{02\_minimal\_model\_sigma.gp} --- over $4582$ fibers ($m\le150$):
all reduction multiplicative, uniform $v_2$ minimalization gap $12$, $v_3$ gap $0$,
$\Delta_{\min}=\Delta_0/2^{12}$, the odd-prime Tate-index formula, and
$N=\rad(\Delta_{\min})$; the exact $\log|\Delta_{\min}|$ identity and the chain
\eqref{eq:chain} with zero violations over $50{,}765$ fibers ($m\le500$);
$\sigma_{\max}=4.6139648$ at $(256,121)$.
\item \texttt{03\_sigma\_density.gp} --- large-sample statistics ($m\le800$,
$129{,}870$ fibers): $\sigma_{\max},\sigma_{\min},\sigma_{\mathrm{mean}}$, and the
density of $\{\sigma\le\sigma_0\}$ compared with the squarefree-locus density at the
three cutoffs $m\le150$, $m\le400$, and $m\le700$.
\item \texttt{04\_sieve\_locus.gp} --- separability of the degree-$8$ form
$ab(a^2-b^2)$, its factorization over $\Q$ into the six forms $m,n,m-n,m+n,F_5,F_6$
(the quartic being reducible over $\Q$ as $F_5F_6$) and the further linear splitting
over $\Q(\sqrt2)$; the containment of $\{\sigma>4\}$ in the union of the six
powerful-value loci (with witnesses driven by both quadratic and linear factors,
e.g.\ $(125,44)$); and the powerful-part exponent and $H^{-1}$ decay of
$\{\sigma>4\}$.
\item \texttt{05\_no\_absolute\_c.gp} --- the all-multiplicative census, the
verified canonical heights $\hath$ (with $\hath(2P)=4\hath(P)$), and the ratio
$\hath/\log|\Delta_{\min}|$ across $\sigma\in[2.75,4.61]$ including the Szpiro-record
fiber $(256,121)$.
\end{itemize}

\section*{Acknowledgements}
The author thanks the maintainers of PARI/GP.

\begin{thebibliography}{99}

\bibitem{bhargava}
M.~Bhargava,
\emph{The geometric sieve and the density of squarefree values of invariant
polynomials},
preprint, arXiv:1402.0031 (2014).

\bibitem{browning}
T.~D.~Browning,
\emph{Power-free values of polynomials},
Arch. Math. (Basel) \textbf{96} (2011), no.~2, 139--150;
building on the determinant method of D.~R.~Heath-Brown,
\emph{The density of rational points on curves and surfaces},
Ann. of Math. (2) \textbf{155} (2002), no.~2, 553--598.

\bibitem{chan}
S.~Chan,
\emph{Almost all elliptic curves with prescribed torsion have Szpiro ratio close
to the expected value},
preprint, arXiv:2407.13850 (2024).

\bibitem{greaves}
G.~Greaves,
\emph{Power-free values of binary forms},
Quart. J. Math. Oxford Ser. (2) \textbf{43} (1992), no.~1, 45--65.

\bibitem{naskrecki}
B.~Naskr\k{e}cki,
\emph{Mordell--Weil ranks of families of elliptic curves associated to
Pythagorean triples},
Acta Arith. \textbf{160} (2013), no.~2, 159--183; arXiv:1210.6933.

\bibitem{pari}
The PARI~Group,
\emph{PARI/GP version 2.15.4},
Univ. Bordeaux, 2023,
\url{https://pari.math.u-bordeaux.fr/}.

\bibitem{peschmann}
R.~Peschmann,
\emph{Quartic reductions and elliptic obstructions for perfect Euler bricks},
preprint, arXiv:2604.09328 (2026).

\bibitem{silverman-atec}
J.~H.~Silverman,
\emph{Advanced Topics in the Arithmetic of Elliptic Curves},
Graduate Texts in Mathematics, vol.~151, Springer, New York, 1994.

\bibitem{stewartyu}
C.~L.~Stewart and K.~Yu,
\emph{On the $abc$ conjecture, II},
Duke Math. J. \textbf{108} (2001), no.~1, 169--181.

\bibitem{yoshida}
T.~Yoshida,
\emph{The relationship between face cuboids and elliptic curves},
preprint, arXiv:2407.09825 (2024).

\end{thebibliography}

\end{document}
